\chapter{Following a store into memory}
\MARSChapterHeading{Following a store into memory}
\label{revision:lab-addresses-and-memory}

This lab accompanies the main book's addresses and memory chapter. Use its
complete \texttt{memory-pixel.asm} source, also supplied in
\texttt{revision/labs/files}. Select MIPS64 Release 6 and the Default memory
layout. The written sequence uses ordinary Run and Debug; this third revised
unit is not yet a native interactive MARS Learning lesson.

The program changes memory rather than printing a result. Its display cell is a
useful visual check, while the byte values let you state exactly what happened.
Work from predictions about addresses and bytes before interpreting a color.

\section{Separate the address from the color}

Write down what the first two instructions under \texttt{main} will put into
\texttt{\$t0} and \texttt{\$t1}. The first value locates \texttt{pixel}; the second
is the color to store there. In this layout, the address is
\texttt{0x10010000} and the color is \texttt{0x00336699}. Both fit in registers,
but their roles differ.

Break at \texttt{store\_color}. At this boundary the address and color have been
prepared, but memory still contains the declared zero. Check both source registers
and the destination word. Predict what changes after one step. A store changes
memory; it does not replace the address register with the color or advance that
register to the next object.

Step once and inspect four bytes beginning at \texttt{pixel}. In increasing
address order, expect \texttt{99 66 33 00}. Explain how that byte sequence agrees
with a word view showing \texttt{00336699}. A disagreement between those two
written forms is not automatically a simulator error.

\section{Track copies across a later write}

Continue to \texttt{change\_blue}. The earlier loads have already completed.
Record \texttt{\$t2}, \texttt{\$t3} and \texttt{\$t4}. The word snapshot is
\texttt{0x00336699}; the unsigned byte is 153; the signed byte is $-103$.
Both byte loads read the same location. Their extension rules explain why their
register values differ.

Step over the byte store. Inspect memory before executing the following load.
Only its lowest-addressed byte should have changed. The three earlier register
values should be unchanged. Explain why a later store cannot retroactively alter
a value that was copied into a register by an earlier load.

Step over the final word loads and compare the earlier word in \texttt{\$t2}
with the current word in \texttt{\$t6}. Check the neighbor in \texttt{\$t7}.
For every observation, name the location and the execution boundary. ``The color
is unchanged'' is ambiguous if it refers to an old register copy rather than to
current memory.

\section{Use the display as a second observation}

The source header connects Bitmap Display to static data and requests two logical
columns and two rows. The pixel is the upper-left cell. Inspect the connected
Debug session's display before and after the byte store. Its red and green
components should remain 51 and 102; blue changes from 153 to 255.

The program may run too quickly for you to distinguish these colors during an
uninterrupted Run. A breakpoint establishes the intermediate boundary. Inspect
the session that actually executed your current source; a previous run's frame
does not establish the new program's result. A displayed color also cannot tell
you whether an unrelated location was overwritten. Keep the memory checks.

\section{Distinguish two faulty stores}

In a copy, change the byte store to use offset one. Before running, compute the
effective address by adding one byte to the base. Predict the word
\texttt{0x0033FF99} and its four bytes. Run and compare. This error changes green
because the offset selects a different byte; the store still writes only one byte.

Restore offset zero and change \texttt{sb} to \texttt{sw}. With 255 in the source,
the four stored bytes become \texttt{FF 00 00 00}. The starting address is now
correct, but the width is wrong. Explain why the two faulty programs both exit
normally even though their results differ from the contract.

As a separate experiment, use byte offset four with \texttt{sb}. Predict which
object it reaches. The address is the first byte of \texttt{neighbor}, so the
pixel remains at its original stored color while the guard changes. Restore the
original program afterward. The offset is measured in bytes for every one of
these instructions; it never counts logical pixels automatically.

\section{Account for every touched byte}

The pixel occupies four consecutive bytes. A word store at offset zero fits that
object. A word store beginning at offset one spans its final three bytes and the
first byte of its neighbor. This Release 6 target supports that unaligned access,
so the program need not stop with an alignment error. Predict the affected bytes
rather than treating the lack of an exception as a bounds check.

An instruction's legal address and the assignment's permitted destination are
different questions. A store can reach writable memory and still damage an object
the program was supposed to preserve. Record both the requested pixel change and
the unchanged neighboring word when you justify your result.

\section{Work from a new component}

Implement the main book's red-component variation in a separate copy. Choose the
byte offset from the little-endian layout rather than from the component's visual
position in the written word. Read the modified component back with \texttt{lbu}
and inspect the complete word as well. Confirm that the other components and guard
are unchanged.

For the memory-copy variation, draw the source, destination and guard as separate
four-byte objects. A register used temporarily during the copy is another
location, not the destination object itself. Trace the load from the source and
the store to the destination, then change the source pixel. Explain why the copied
word remains independent. Use that trace to diagnose any version in which both
objects unexpectedly change or the destination was never written.
